% ═══════════════════════════════════════════════════════════════════════════════
%  CAPÍTULO — Template method
% ═══════════════════════════════════════════════════════════════════════════════

\chapter{Template Method}

% ─── Situación Problema ──────────────────────────────────────────────────────
\section{Situación Problema}

El proceso de publicación de un nuevo alojamiento en Airbnb sigue siempre
los mismos pasos generales: validar los datos básicos, subir y verificar
las fotografías, configurar el calendario de disponibilidad, definir el
precio base y finalmente publicar el anuncio. Sin embargo, algunos pasos
varían según el tipo de alojamiento: publicar una \emph{experiencia}
requiere validar las credenciales del anfitrión-guía; publicar una
\emph{propiedad de lujo} activa una revisión adicional por parte de
un equipo especializado.

Si cada tipo de alojamiento reimplementa el flujo completo de publicación,
se duplica la lógica común y aumenta el riesgo de que los tipos difieran
en pasos que deberían ser idénticos, como la verificación de fotografías
o la publicación final.

% ─── Solución ────────────────────────────────────────────────────────────────
\section{Solución}

El patrón Template Method define en la clase base
\texttt{ProcesoPublicacion} el esqueleto del algoritmo de publicación
como un método \texttt{publicar()} que llama en orden a los pasos del
proceso. Los pasos invariantes —validación de datos, verificación de
fotos, publicación final— se implementan directamente en la clase base;
los pasos específicos de cada tipo se declaran como métodos abstractos
que las subclases concretas (\texttt{PublicacionExperiencia},
\texttt{PublicacionPropiedadLujo}) implementan a su manera.

La secuencia global está protegida en un único lugar y no puede ser
alterada por las subclases, garantizando consistencia. Solo los puntos
de variación reales quedan abiertos a especialización, evitando la
duplicación del esqueleto del proceso.

% ─── Diagrama de clases ─────────────────────────────────────────────────────
\begin{figure}[H]
    \centering
    % \includegraphics[width=\textwidth]{imagenes/abstract_factory_diagrama.png}
    \caption{Diagrama de clases del patrón Abstract Factory}
\end{figure}


% ─── Roles de las Clases ────────────────────────────────────────────────────
\section{Roles de las Clases}

% Explicar la responsabilidad de cada clase/interfaz

\begin{itemize}
    \item \textbf{AbstractFactory:}
    \item \textbf{ConcreteFactory:}
    \item \textbf{AbstractProduct:}
    \item \textbf{ConcreteProduct:}
    \item \textbf{Client:}
\end{itemize}


% ─── Main ───────────────────────────────────────────────────────────────────
\section{Main}

% Explicar qué realiza el programa principal

\begin{lstlisting}[language=Java, caption=NombreClase]
 
\end{lstlisting}


% ─── Salida del Main ────────────────────────────────────────────────────────
\section{Salida del Main}

\begin{lstlisting}[language=Java, caption=NombreClase]
 
\end{lstlisting}


% ─── Clases ─────────────────────────────────────────────────────────────────
\section{Clases}

% ─── Clase 1 ────────────────────────────────────────────────────────────────
\subsection{NombreClase}

\begin{lstlisting}[language=Java, caption=NombreClase]
 
\end{lstlisting}


% ─── Clase 2 ────────────────────────────────────────────────────────────────
\subsection{NombreClase}

\begin{lstlisting}[language=Java, caption=NombreClase]
 
\end{lstlisting}


% ─── Clase 3 ────────────────────────────────────────────────────────────────
\subsection{NombreClase}

\begin{lstlisting}[language=Java, caption=NombreClase]
 
\end{lstlisting}


% ─── Conclusiones ───────────────────────────────────────────────────────────
\section{Conclusiones}

% Explicar:
% - Ventajas del patrón
% - Cuándo usarlo
% - Beneficios obtenidos
% - Posibles desventajas
% - Relación con principios SOLID